\section{Estado del Arte}\label{sec:estado_arte}

El desarrollo de soluciones modernas para la operación eficiente, el diagnóstico predictivo y la seguridad humana en accionamientos industriales de corriente alterna exige comprender la convergencia tecnológica ocurrida entre 2019 y 2026. A continuación se delimita la problemática industrial que motiva este proyecto, se revisan críticamente 25 investigaciones científicas recientes organizadas en tres ejes complementarios, se sintetiza una comparativa técnica y se establece cómo la plataforma propuesta cierra la brecha identificada.

\subsection{Identificación de la Problemática Industrial}\label{subsec:problematica_ea}

En plantas agroindustriales de molienda de granos, los motores trifásicos de inducción jaula de ardilla constituyen el núcleo motriz de la cadena productiva y concentran una fracción mayoritaria del consumo eléctrico operativo. Sin embargo, la literatura técnica reciente evidencia cuatro problemas persistentes que comprometen simultáneamente la competitividad energética, la disponibilidad mecánica y la integridad física del personal:

\begin{enumerate}\setlength{\itemsep}{2pt}\setlength{\topsep}{2pt}
\item \textbf{Pérdidas ocultas de eficiencia y sobrecostos del mantenimiento reactivo:} las fallas incipientes en devanados, barras rotóricas, rodamientos o alimentación degradan progresivamente el rendimiento energético y elevan las pérdidas por efecto Joule mucho antes de provocar un disparo catastrófico~\cite{Garcia2019}. Cuando la gestión de activos depende de esquemas correctivos o preventivos rígidos basados únicamente en calendario, las paradas no programadas y los daños secundarios multiplican los costos de operación y mantenimiento~\cite{Maintenance2023}.
\item \textbf{Ineficiencia en arranques directos (DOL) y regímenes de carga parcial:} numerosos procesos de trituración operan a frecuencia fija de red (\qty{60}{\hertz}) sin modulación adaptativa de velocidad ni flujo magnético. Esta práctica genera sobreconsumo energético en etapas de flujo másico reducido, picos severos de corriente de rotor bloqueado y pulsaciones de par electromagnético que fatigan la transmisión mecánica~\cite{Resa2019,Shabestari2019}.
\item \textbf{Escasez de datos industriales etiquetados de falla:} aunque la inteligencia artificial (IA) promete diagnósticos tempranos, las plantas reales rara vez permiten llevar motores críticos hasta la rotura para recopilar bases de datos balanceadas. Esta carencia limita la generalización de modelos supervisados al pasar del laboratorio al entorno fabril~\cite{Soualhi2022,AI2026}.
\item \textbf{Desconexión entre telemetría de condición y capacitación segura del operario:} las arquitecturas convencionales de Industria~4.0 priorizan tableros analíticos de variables eléctricas, pero descuidan la dimensión humana exigida por la Industria~5.0~\cite{Breque2021,Xu2021}. En consecuencia, los técnicos carecen de entornos inmersivos donde practicar maniobras de bloqueo y etiquetado (LOTO), uso reglamentario de elementos de protección personal (EPP) y respuesta ante fallas eléctricas reales sin exponerse a choques ni arcos eléctricos~\cite{Stefan2023,Stefan2024}.
\end{enumerate}

\subsection{Eje 1: Eficiencia Energética y Control en Accionamientos AC (2019--2024)}\label{subsec:eje1_eficiencia}

La evolución investigativa en eficiencia energética de máquinas de inducción demuestra una transición clara desde el análisis aislado de pérdidas hacia estrategias de control adaptativo integradas con diagnóstico en servicio. En 2019, García \emph{et~al.}~\cite{Garcia2019} evaluaron cuantitativamente el impacto de distintas condiciones de falla incipiente sobre la eficiencia de motores de inducción en operación, comprobando que las asimetrías electromecánicas incrementan el deslizamiento, distorsionan la corriente estatórica y penalizan de forma sostenida el consumo eléctrico. Ese mismo año, Resa \emph{et~al.}~\cite{Resa2019} propusieron una estrategia de reducción del consumo energético mediante referencias variables conjuntas de velocidad y flujo magnético, evidenciando ahorros significativos frente a la operación a flujo nominal constante bajo perfiles de carga fluctuantes. Complementariamente, Shabestari y Mehrizi-Sani~\cite{Shabestari2019} abordaron el control de limitación de corriente y la atenuación del rizado de par electromagnético en motores de inducción, reduciendo tanto los esfuerzos térmicos en transitorios de arranque como la fatiga mecánica en ejes y acoples.

Durante 2020, Goun \emph{et~al.}~\cite{Goun2020} avanzaron hacia la implementación práctica mediante la determinación experimental de parámetros subóptimos para el control energéticamente eficiente de motores de inducción, demostrando que leyes de control calibradas experimentalmente capturan gran parte del potencial de ahorro sin exigir modelos analíticos excesivamente costosos en tiempo real. Con la adopción masiva de variadores de frecuencia (VFD), la complejidad espectral de los accionamientos aumentó. En 2023, Yakhni \emph{et~al.}~\cite{Yakhni2023} revisaron las técnicas de análisis transitorio de la firma de corriente del motor (MCSA) en accionamientos de velocidad variable, señalando que el contenido armónico de los inversores y las trayectorias no estacionarias de velocidad invalidan las transformadas estacionarias clásicas y exigen procesamiento transitorio o multivariable. En paralelo, Zhang \emph{et~al.}~\cite{Review2023} presentaron una revisión integral de las estrategias de diagnóstico de fallas en trenes motrices electrificados, destacando la interacción crítica entre los lazos cerrados de control vectorial o de par directo (DTC) y la observabilidad de anomalías eléctricas y mecánicas.

En 2024, Adamou y Alaoui~\cite{Adamou2024} consolidaron la convergencia entre gestión energética y monitoreo de condición mediante un enfoque de diagnóstico centrado en la eficiencia para gemelos digitales de motores de inducción en servicio. Su estudio sobre barras rotóricas rotas demostró que vigilar desviaciones de eficiencia en tiempo real proporciona un indicador complementario de salud útil tanto para mantenimiento como para auditoría energética. No obstante, los estudios de este eje rara vez vinculan el desempeño del accionamiento con indicadores productivos específicos por tonelada procesada ni traducen el ahorro eléctrico en mitigación verificable de huella de carbono según el factor de emisión de la red nacional.

\subsection{Eje 2: Diagnóstico Predictivo con IA, Fusión Sensorial y Arquitectura Edge/Fog (2020--2026)}\label{subsec:eje2_ia_edge}

El segundo eje del estado del arte refleja la maduración del mantenimiento predictivo asistido por aprendizaje automático y arquitecturas distribuidas de adquisición. En 2020, Gangsar y Tiwari~\cite{Gangsar2020} demostraron la eficacia de las máquinas de vectores de soporte (SVM) para el diagnóstico multiclase de fallas eléctricas y mecánicas en motores de inducción mediante la fusión de señales de vibración y corriente estatórica, confirmando que la combinación multisensorial supera las limitaciones de detección de cada magnitud por separado bajo cargas variables. En 2021, Glowacz~\cite{Glowacz2021} exploró el monitoreo no invasivo basado en señales acústicas y redes neuronales artificiales (ANN), validando la capacidad de clasificadores neuronales ligeros para discriminar estados anómalos incipientes con bajo costo computacional.

En 2022 se produjeron cuatro aportes estructurantes. Soualhi \emph{et~al.}~\cite{Soualhi2022} publicaron una revisión exhaustiva sobre el uso de la IA en la evaluación de salud de motores de inducción, abarcando extracción de características, clasificación de fallas, índices de salud y pronóstico de vida útil remanente (RUL), e identificando la interpretabilidad clínica del diagnóstico y la robustez ante cambios de régimen como desafíos abiertos. En el ámbito del aprendizaje profundo, Ribeiro Junior \emph{et~al.}~\cite{Ribeiro2022} implementaron redes neuronales convolucionales unidimensionales (1D-CNN) sobre canales múltiples de vibración para automatizar la extracción espectral y elevar la precisión multiclase. Para viabilizar el despliegue fabril continuo, Li \emph{et~al.}~\cite{EdgeCloud2022} propusieron una arquitectura IIoT \emph{Edge-to-Cloud} con sensores inteligentes ubicuos que separa el preprocesamiento local de baja latencia de la analítica histórica en la nube, mientras Cinar~\cite{SensorFusion2022} combinó fusión sensorial con aprendizaje por transferencia (\emph{transfer learning}) para mejorar la generalización diagnóstica cuando el volumen de muestras etiquetadas es restringido.

Durante 2023, la integración entre IoT y aprendizaje profundo se consolidó en múltiples direcciones. Moleda \emph{et~al.}~\cite{Maintenance2023} sistematizaron la transición desde el mantenimiento correctivo hasta el predictivo en la industria eléctrica y energética, cuantificando los beneficios en confiabilidad al sustituir inspecciones estáticas por analítica continua. Tran \emph{et~al.}~\cite{Tran2023} desarrollaron un esquema robusto de reconocimiento y corrección de fallas en motores de inducción acoplando telemetría IoT con redes profundas, mientras Kumar \emph{et~al.}~\cite{Kumar2023} ratificaron la superioridad de las arquitecturas de fusión multisensorial con aprendizaje profundo para clasificar simultáneamente fallas eléctricas y mecánicas en condiciones operativas exigentes.

En el periodo 2025--2026, la frontera investigativa ha atacado directamente tres cuellos de botella críticos: latencia de cómputo, degradación térmica y escasez de datos. En 2025, D'Agostino \emph{et~al.}~\cite{Fog2025} formularon una arquitectura escalable de computación en la niebla (\emph{fog computing}) para mantenimiento predictivo industrial personalizable, demostrando que ejecutar inferencia cerca del activo elimina cuellos de botella de conectividad y refuerza la autonomía operativa. Ese mismo año, Zhang \emph{et~al.}~\cite{Thermal2025} desarrollaron un modelado térmico predictivo habilitado por gemelo digital para el monitoreo continuo de la temperatura del estator en motores de inducción, permitiendo anticipar puntos calientes y envejecimiento del aislamiento antes del daño dieléctrico irreversible. Finalmente, en 2026, Struwig \emph{et~al.}~\cite{AI2026} abordaron el monitoreo de salud basado en IA para motores de inducción Clase~I en entornos con escasez de datos, estableciendo una metodología rigurosa que enlaza la generación sintética de líneas base y escenarios de falla con el despliegue industrial real. A pesar de estos avances, la mayoría de las propuestas concluye en una etiqueta de alarma en pantalla o en plataformas web desconectadas del proceso físico, sin incorporar validación cruzada embebida de dos familias de clasificadores ni listas de chequeo prescriptivas de mantenimiento centrado en confiabilidad (RCM).

\subsection{Eje 3: Gemelos Digitales, Industria 5.0 y Juegos Serios para Seguridad Eléctrica (2021--2024)}\label{subsec:eje3_i50_seguridad}

El tercer eje articula la transformación de los sistemas ciberfísicos hacia plataformas centradas en el ser humano. En 2021, el informe de referencia de la Comisión Europea elaborado por Breque \emph{et~al.}~\cite{Breque2021} definió los tres pilares rectores de la Industria~5.0: orientación al ser humano, sostenibilidad ambiental y resiliencia industrial. En concordancia, Xu \emph{et~al.}~\cite{Xu2021} analizaron el tránsito conceptual de la Industria~4.0 hacia la Industria~5.0, aclarando que esta última no reemplaza los habilitadores digitales previos (IoT, IA y gemelos digitales), sino que los subordina a la seguridad, el empoderamiento cognitivo del operario y la descarbonización de los procesos productivos.

En 2023, la convergencia entre gemelos digitales y formación segura cobró impulso en dos líneas hasta entonces paralelas. Por un lado, Singh \emph{et~al.}~\cite{Singh2023} diseñaron un sistema de mantenimiento predictivo inteligente soportado por un gemelo digital para máquinas eléctricas AC industriales, demostrando que la sincronización entre modelos electromecánicos y analítica predictiva agiliza la toma de decisiones en planta. Por otro lado, Stefan \emph{et~al.}~\cite{Stefan2023} publicaron una revisión sistemática sobre la efectividad de la realidad virtual (VR) y los entornos 3D interactivos en el entrenamiento relacionado con seguridad laboral, concluyendo que la práctica experiencial en escenarios virtuales mejora la identificación de peligros y la retención procedimental frente a la instrucción teórica pasiva.

En 2024, ambas vertientes mostraron madurez adicional, aunque aún fragmentada. Fera y Spandonidis~\cite{Maritime2024} integraron el diagnóstico de fallas basado en IA dentro de un marco unificado de gemelo digital para sistemas eléctricos de potencia marítimos, evidenciando el valor de acoplar modelos de planta con soporte inteligente a decisiones en sistemas críticos. Simultáneamente, Stefan \emph{et~al.}~\cite{Stefan2024} evaluaron empíricamente la efectividad de los entornos virtuales específicamente en capacitación de seguridad eléctrica ---incluyendo procedimientos de aislamiento, reconocimiento de peligros energizados y verificación de protocolos---, reportando mejoras verificables en aprendizaje, aceptación del participante y eficiencia del tiempo de entrenamiento sin exposición real al riesgo eléctrico. Sin embargo, los gemelos digitales de máquinas eléctricas~\cite{Singh2023,Maritime2024} no integran capacitación física en primera persona sobre EPP y reglas de oro LOTO, mientras que los simuladores de seguridad eléctrica~\cite{Stefan2023,Stefan2024} prescinden de física electromecánica continua, optimización energética VFD y clasificación IA en tiempo real.

\begin{table*}[!t]
\begingroup\setlength{\tabcolsep}{3.5pt}\renewcommand{\arraystretch}{0.95}
\caption{Comparativa técnica de siete trabajos emblemáticos del estado del arte (2019--2026) frente a la solución propuesta.}
\label{tab:estado_arte}
\centering\footnotesize
\begin{tabularx}{\textwidth}{@{} R{2.3cm} c R{3.2cm} R{3.5cm} Y @{}}
\toprule
\textbf{Estudio} & \textbf{Año} & \textbf{Técnica IA / Control} & \textbf{Gemelo 3D / Seguridad Humana} & \textbf{Optimización Energética y Brecha Resuelta} \\
\midrule
García \emph{et~al.}~\cite{Garcia2019} & 2019 & Evaluación analítico-experimental de pérdidas bajo fallas & Sin entorno 3D interactivo ni entrenamiento humano & Cuantifica la caída de eficiencia por fallas incipientes; no incorpora control VFD, clasificación IA ni protocolos de intervención segura. \\
\addlinespace[2pt]
Gangsar y Tiwari~\cite{Gangsar2020} & 2020 & SVM multiclase fusionando vibración y corriente estatórica & Sin gemelo digital 3D ni formación en seguridad eléctrica & Logra alta precisión diagnóstica a velocidad variable; carece de pronóstico RUL, análisis de consumo específico y soporte RCM en planta. \\
\addlinespace[2pt]
Li \emph{et~al.}~\cite{EdgeCloud2022} & 2022 & IIoT \emph{Edge-to-Cloud} con sensores inteligentes ubicuos & Supervisión 2D remota; sin entrenamiento EPP/LOTO & Reduce latencia de monitoreo industrial; depende de infraestructura en la nube y omite optimización energética de accionamientos AC. \\
\addlinespace[2pt]
Singh \emph{et~al.}~\cite{Singh2023} & 2023 & Gemelo digital con mantenimiento predictivo inteligente & Réplica virtual de máquina AC; sin interacción física de operario & Une monitoreo de condición y diagnóstico en máquinas AC; no contempla seguridad normativa (RETIE/LOTO) ni ahorro energético con huella de $\mathrm{CO_2}$. \\
\addlinespace[2pt]
Stefan \emph{et~al.}~\cite{Stefan2024} & 2024 & Simulación interactiva enfocada en instrucción procedimental & Entorno virtual para capacitación en seguridad eléctrica & Valida ganancia pedagógica sin riesgo físico; carece de modelo multifísico del motor, telemetría SCADA, IA predictiva y control VFD. \\
\addlinespace[2pt]
Zhang \emph{et~al.}~\cite{Thermal2025} & 2025 & Gemelo digital térmico predictivo de devanado estatórico & Modelo virtual térmico; sin inmersión 3D ni interacción segura & Anticipa sobrecalentamiento y envejecimiento del aislamiento; no abarca fallas mecánicas/eléctricas múltiples ni gestión energética VFD. \\
\addlinespace[2pt]
Struwig \emph{et~al.}~\cite{AI2026} & 2026 & IA con generación sintética en escasez de datos (Clase~I) & Sin entorno 3D interactivo ni evaluación de seguridad humana & Resuelve el cuello de botella de datos de falla en motores industriales; no integra control PID/VFD, huella ambiental ni Industria~5.0. \\
\midrule
\textbf{Este trabajo (AgroMolinos 3D)} & \textbf{2026} & \textbf{Simulink DTC + Kloss \qty[mode=text]{10}{HP}, MLP + SVM RBF nativos, VFD $V/f$ + PID} & \textbf{Planta 3D Godot~4.7, 5 fallas diegéticas, EPP 6/6, 5 Reglas LOTO y rúbrica RETIE} & \textbf{Cierra la brecha integral: ahorro de \qty[mode=text]{28.38}{\percent} (\qty[mode=text]{7.52}{kWh/ton} vs.\ \qty[mode=text]{10.50}{kWh/ton}), \qty[mode=text]{488.7}{\kilo\gram} de $\mathbf{CO_2e}$/año evitados (UPME) y Tablet SCADA embebida sin web externa.} \\
\bottomrule
\end{tabularx}
\endgroup
\end{table*}

\subsection{Síntesis de la Brecha y Solución Propuesta}\label{subsec:sintesis_brecha}

La Tabla~\ref{tab:estado_arte} contrasta siete contribuciones representativas del periodo 2019--2026 frente al sistema desarrollado en este proyecto. El análisis transversal revela que, si bien existen avances maduros por separado en eficiencia energética bajo fallas~\cite{Garcia2019,Resa2019,Adamou2024}, diagnóstico con IA en el borde y gemelos térmicos~\cite{Gangsar2020,EdgeCloud2022,Fog2025,Thermal2025,AI2026} y entrenamiento virtual en seguridad eléctrica~\cite{Stefan2023,Stefan2024}, persiste una brecha estructural: \emph{ninguna solución revisada unifica en un único entorno autónomo la física electromecánica de un accionamiento industrial, la inferencia Edge-AI dual con prescripción RCM, la optimización energética con contabilidad ambiental y la formación diegética en seguridad eléctrica bajo normativa colombiana}.

Para cerrar integralmente esta brecha y materializar la transición de Industria~4.0 a Industria~5.0 en una planta molino de granos (\emph{AgroMolinos del Valle S.A.}), la solución propuesta articula cuatro innovaciones acopladas:

\begin{itemize}\setlength{\itemsep}{2pt}\setlength{\topsep}{2pt}
\item \textbf{Modelado multifísico coherente (Simulink DTC + Godot Kloss):} se caracteriza el motor trifásico de inducción jaula de ardilla de \qty{10}{HP} (\qty{7.46}{\kilo\watt}), \qty{460}{\volt}, \qty{60}{\hertz} y 4~polos ($N_r = \qty{1750}{RPM}$, $I_{\mathrm{FLA}} = \qty{13.5}{\ampere}$, $I_{\mathrm{LRA}} = \qty{81.0}{\ampere}$, $T_{\max} = \qty{101.75}{\newton\meter}$) mediante simulación detallada de control directo de par en MATLAB/Simulink y se traslada su ecuación electromecánica de Kloss, su dinámica térmica de aislamiento Clase~F (\qty{155}{\celsius}) y su respuesta vibratoria según ISO~10816-3/20816-3~\cite{ISO10816,ISO20816} al motor físico en tiempo real a \qty{10}{\hertz} en Godot~4.7. Esta correspondencia física permite generar conjuntos de datos sintéticos físicamente consistentes que superan la escasez industrial de fallas etiquetadas~\cite{AI2026,Thermal2025}.
\item \textbf{Diagnóstico dual embebido (MLP + SVM RBF) y analítica RCM sin nube externa:} sobre un vector de ocho variables True-RMS ($I_{\mathrm{RMS}}$, factor de desbalance $\mathrm{IUF}$, distorsión armónica $\mathrm{THD}$, $V_{\mathrm{RMS}}$, $T_w$, $dT/dt$, $\mathrm{Vib}_{\mathrm{RMS}}$ y $\mathrm{Vib}_{\mathrm{peak}}$), se ejecutan nativamente en el simulador una red perceptrón multicapa MLP ($8 \to 16 \to 8 \to 4$, exactitud de \qty{99.17}{\percent}) y una máquina de vectores de soporte SVM con núcleo RBF (exactitud de \qty{98.75}{\percent}). Ambos clasificadores se contrastan en tiempo real para emitir el índice de salud del motor (MHI), la vida útil remanente (RUL) por degradación térmica de Arrhenius y listas verificables de mantenimiento RCM en una Tablet SCADA industrial nativa (\texttt{[TAB]}), eliminando la dependencia de servidores o navegadores web externos~\cite{Fog2025,Soualhi2022}.
\item \textbf{Optimización energética VFD/PID y huella de carbono UPME:} el controlador integrado permite conmutar en operación entre arranque directo DOL (\qty{10.50}{kWh/ton}, pico de \qty{81.0}{\ampere}) y variador escalar $V/f$ con lazo PID proporcional al flujo másico de grano (\qty{7.52}{kWh/ton}, corriente de arranque menor a \qty{17.5}{\ampere}). Aprovechando la ley cúbica de afinidad en carga parcial (validada a \qty{42}{\hertz}), se alcanza un ahorro energético específico de \qty{28.38}{\percent} y una reducción anual estimada de \qty{488.7}{\kilo\gram} de $\mathrm{CO_2e}$ calculada con el factor oficial del Sistema Interconectado Nacional de Colombia publicado por la UPME (\qty{164.38}{g/kWh} de $\mathrm{CO_2}$)~\cite{UPME2024}.
\item \textbf{Entrenamiento inmersivo diegético en seguridad RETIE y LOTO:} en lugar de inyectar anomalías mediante botones bidimensionales artificiales, el estudiante induce y repara físicamente cinco fallas reales en el entorno 3D (corte de fase L3, atasco en tolva con sacos de \qty{50}{\kilo\gram}, caída de tensión a \qty{340}{\volt} en el CCM, obstrucción de ventilación trasera y soltura de pernos de bancada) mediante herramientas físicas dedicadas. El simulador impone enclavamiento obligatorio de EPP 6/6, secuencia estricta de las 5 Reglas de Oro LOTO, perímetro de arco eléctrico de \qty{1.0}{\meter}, curvas fisiológicas IEC~60479-1 y una rúbrica continua auditable de \num{0.0} a \num{5.0} alineada con el RETIE y OSHA~1910.147~\cite{RETIE2026,OSHA1910,IEC60479}.
\end{itemize}
